% TL;DR: Inhalt der Projektkarte (U-Bahn-Plan April–Nov. 2026) ohne Dokumentrahmen; eingebunden von karte.tex (build.sh) und direkt von der Ausarbeitung (Kapitel Projektverlauf). Quelle der Zahlen: docs/ergebnisse/diagnose-chronik.md.
%
% Voraussetzung: karte_stil.tex ist geladen. Variante über \def\kartevariante{voll|kurz} vor \input:
%   voll = alle Stationen mit Detailzeile (A4, Poster, README, Ausarbeitung)
%   kurz = nur Wendepunkte und markierte Stationen, ohne Detailzeile (Folie)
% Alle Definitionen hier bleiben lokal, wenn die Datei in einer Gruppe eingebunden wird.
%
\providecommand\kartevariante{voll}%
\newif\ifkurz%
\def\tmpkurz{kurz}%
\ifx\kartevariante\tmpkurz \kurztrue \fi%
%
% Im Kurzmodus: Beschriftung nur für Wendepunkte (*) und Schlüsselstationen (!)
\ifkurz%
  \RenewDocumentCommand\St{s t! m m m m m m}{%
    \IfBooleanTF{#1}{\Wende{#4,#5}\def\sep{\lblsepW}}{\def\sep{\lblsep}}%
    \Mk{#3}{#4,#5}%
    \IfBooleanTF{#1}{\node[lbl, #6=\sep] at (#4,#5) {{\bfseries #7}};}{%
      \IfBooleanT{#2}{\node[lbl, #6=\sep] at (#4,#5) {{\bfseries #7}};}}}%
\else%
  \RenewDocumentCommand\St{s t! m m m m m m}{%
    \IfBooleanTF{#1}{\Wende{#4,#5}\def\sep{\lblsepW}}{\def\sep{\lblsep}}%
    \Mk{#3}{#4,#5}%
    \node[lbl, #6=\sep] at (#4,#5) {{\bfseries #7}\ifx\relax#8\relax\else\\#8\fi};}%
\fi%
\newcommand\Detail[1]{\ifkurz\else#1\fi}%
\newcommand\fsHub{\fsLbl}% Schrift der Balken im Kasten
\ifkurz % Folie: größere Schrift, weniger Text
  \renewcommand\fsLbl{\fontsize{8.4}{9.4}\selectfont\sffamily}%
  \renewcommand\fsPhase{\fontsize{8.4}{9.6}\selectfont\sffamily}%
  \renewcommand\fsHub{\fontsize{7.2}{8.2}\selectfont\sffamily}%
  \tikzset{badge/.append style={font=\fontsize{8}{8.8}\selectfont\sffamily\bfseries, minimum width=19mm, inner xsep=1mm}}%
\fi%
%
\begin{tikzpicture}[x=1cm, y=1cm]

% ── Lagen der Linien (y) ──────────────────────────────────────────────
\def\yI{0}      % Infrastruktur
\def\yT{-2.0}   % Training N1.6
\def\yN{-3.35}  % N1.7
\def\yC{-5.1}   % Co-Training / Daten
\def\yS{-7.1}   % Simulation / Szene
\def\yD{-9.6}   % Evaluation / Diagnose
\def\yR{-12.3} % RL
\def\yK{-13.45} % Doku / Abgabe
\def\xH{24.0}   % linke Kante "Stand heute"

% ── Phasenbänder (Zeitachse nicht maßstäblich) ────────────────────────
\begin{scope}[on background layer]
  \foreach \xa/\xb/\cc/\nam/\dat in {
      -2.4/0.0/phaseB/{Eigene Rechner}/{April},
      0.0/4.8/phaseA/{Aufbau \& erste Läufe}/{08.05.–05.06.},
      4.8/6.6/white/{Prüfungs-\\phase}/{Juni–Juli,\\reduziert},
      6.6/11.6/phaseA/{Sim auf Blackwell, Diagnose}/{01.–14.08.},
      11.6/14.8/phaseB/{Geometrieabgleich}/{15.–26.08.},
      14.8/17.0/phaseA/{Synthesedaten \& N1.7}/{27.08.–11.10.},
      17.0/20.4/phaseB/{Nachmessen \& Registrierung}/{12.–14.10.},
      20.4/23.8/phaseA/{Co-Training v3}/{15.–29.10.}} {
    \pgfmathsetlengthmacro\tw{(\xb-\xa)*1cm-2mm}
    \fill[\cc] (\xa,2.7) rectangle (\xb,-14.05);
    \node[font=\fsPhase, align=center, text width=\tw, anchor=north]
      at ({(\xa+\xb)/2},2.62) {{\bfseries\nam}\\[-0.5pt]{\color{black!60}\dat}};
  }
  \fill[pattern=north east lines, pattern color=black!12] (4.8,1.25) rectangle (6.6,-14.05);
  \draw[black!35, line width=0.4pt] (-2.4,1.25) -- (23.8,1.25);
  % Monatsanfänge, innerhalb der Phasen linear umgerechnet (z. B. 01.06. = Tag 24 von 28 in Phase 1)
  \foreach \x/\m in {-0.45/Mai, 4.11/Jun, 5.6/Jul, 6.6/Aug, 15.04/Sep, 16.51/Okt} {
    \draw[black!35, line width=0.4pt] (\x,1.25) -- (\x,1.1);
    \node[font=\fontsize{5.4}{6}\selectfont\sffamily, text=black!50, anchor=north west,
      inner sep=0.4pt] at (\x,1.2) {1.\,\m};
  }
\end{scope}

% ── Linien-Plaketten links ────────────────────────────────────────────
\foreach \y/\cc/\lab in {\yI/lInfra/{Infrastruktur},\yT/lTrain/{Training N1.6},
    \yN/lNsieben/{GR00T N1.7},\yC/lCotr/{Co-Training\\[-1pt]\& Daten},
    \yS/lSim/{Simulation\\[-1pt]\& Szene},\yD/lDiag/{Evaluation\\[-1pt]\& Diagnose},
    \yR/lRL/{Reinforcement\\[-1pt]Learning},\yK/lDoku/{Doku \& Abgabe}}
  \node[badge, fill=\cc] at (-3.65,\y) {\lab};

% ════════════════════════════ LINIEN ══════════════════════════════════
% Infrastruktur
\draw[line, lInfra] (-2.65,\yI) -- (\xH,\yI);
\draw[stub, lInfra] (2.3,\yI) -- (2.85,-0.55) -- (3.45,-0.55); \Prellbock[lInfra]{3.45,-0.55}{0}
\draw[stub, lInfra] (6.0,\yI) -- (6.55,-0.55) -- (7.2,-0.55); \Prellbock[lInfra]{7.2,-0.55}{0}
\draw[stub, lInfra] (-1.35,\yI) -- (-0.85,-0.75) -- (-0.35,-0.75); \Prellbock[lInfra]{-0.35,-0.75}{0}

% Training N1.6
\draw[line, lTrain] (-2.65,\yT) -- (\xH,\yT);
\draw[stub, lTrain] (2.5,\yT) -- (3.05,-2.6) -- (3.7,-2.6); \Prellbock[lTrain]{3.7,-2.6}{0}
% Fehlannahme U-Kurve (14.08. → 12.10.)
\draw[irrweg, lTrain] (11.25,\yT) -- (11.7,-2.48) -- (16.95,-2.48) -- (17.25,\yT);

% N1.7 (zweigt am 19.08. ab)
\draw[line, lNsieben] (12.4,\yT) -- (13.75,\yN) -- (\xH,\yN);

% Co-Training / Daten
\draw[line, lCotr] (11.9,\yC) -- (\xH,\yC);
\draw[stub, lCotr] (12.75,\yC) -- (13.3,-5.65) -- (13.8,-5.65); \Prellbock[lCotr]{13.8,-5.65}{0}

% Simulation / Szene
\draw[line, lSim] (-2.65,\yS) -- (\xH,\yS);
% Fehlannahme Finger-Vorzeichen (≤ 04.06. → 24.08.)
\draw[irrweg, lSim] (3.55,\yS) -- (4.25,-8.3) -- (13.5,-8.3) -- (13.95,\yS);

% Evaluation / Diagnose
\draw[line, lDiag] (-2.65,\yD) -- (\xH,\yD);
% Fehlannahme SigLIP-Abstand (04.06. → 14.10.)
\draw[irrweg, lDiag] (3.3,\yD) -- (3.95,-10.95) -- (18.6,-10.95) -- (19.0,\yD);

% RL
\draw[line, lRL] (4.6,\yR) -- (8.4,\yR);
\draw[stub, lRL] (8.4,\yR) -- (\xH,\yR);

% Doku
\draw[line, lDoku, line width=1.6pt] (-2.65,\yK) -- (\xH,\yK);

% Umsteigekapsel Co-Training v3 (Training, N1.7, Daten)
\draw[black, line width=0.9pt, fill=white, rounded corners=2.4mm]
  (20.95,{\yT+0.25}) rectangle (21.45,{\yC-0.25});

% ════════════════════════════ STATIONEN ═══════════════════════════════
% Infrastruktur
% April: eigene Rechner unter Windows; ab 08.05. nur noch Linux + Docker (x < 0 = April)
\St!{old}{-1.7}{\yI}{above}{Eigene PCs}{Windows + WSL2}
\St{fail}{-0.35}{-0.75}{right}{Training + Sim lokal}{\Detail{zu wenig VRAM,\\Sim lief nicht}}
\St*!{ok}{0.3}{\yI}{above}{Linux + Docker}{Fork, 08.05.}
\St{ok}{1.95}{\yI}{above}{KISSKI-Cluster}{4× A100}
\St{ok}{3.75}{\yI}{above}{Sim auf Miet-GPU}{vast.ai}
\St{fail}{3.45}{-0.55}{right}{Sim auf KISSKI}{\Detail{keine RT-Cores}}
\St{ok}{5.6}{\yI}{above}{Institutsserver}{Blackwell, Juli}
\St{fail}{7.2}{-0.55}{right}{Isaac Sim 5.1}{\Detail{instabil}}
\St*{ok}{7.6}{\yI}{above}{Isaac Sim 6.0}{08.08.}
\St{ok}{12.4}{\yI}{above}{./run.sh-Menüs}{20.08.}
\St{ok}{15.6}{\yI}{above}{Image je Generation}{29.09.}
\St{ok}{18.4}{\yI}{above}{CI + Reviews}{13.10.}
\St{part}{22.4}{\yI}{above}{Watchdog}{Prozess hängt, Ursache offen}

% Training N1.6
\St!{ok}{1.6}{\yT}{above}{T1}{175k Schritte, Encoder fest}
\St{fail}{3.7}{-2.6}{right}{T2: Encoder trainiert}{\Detail{Policy kollabiert}}
\St{ok}{9.65}{\yT}{above}{Split aktiv}{13.08.}
\St*{ok}{11.25}{\yT}{above}{T3}{Vision + Jitter + Split}
\node[lbl, lTrain, font=\fsLbl\itshape] at (15.35,-2.74) {Fehlannahme: U-Kurve, »Optimum @30k«};
\St*{ok}{17.25}{\yT}{above}{Nachmessung T3}{20×750: @44k am besten}
\St{part}{22.7}{\yT}{above}{N1.6 cube4 @44k}{frei gehoben 2/60,\\gestapelt 0/60}

% N1.7
\St!{ok}{14.25}{\yN}{below}{N1.7-Pfad}{19.08.}
\St{ok}{17.25}{\yN}{below}{T4 (N1.7)}{gleichauf, 1,6× schneller}
\St{fail}{19.2}{\yN}{above}{N1.7 in der Sim}{hebt 0/10 (L54/55)}
\St{part}{22.7}{\yN}{above}{N1.7 co-trainiert}{frei gehoben 0/20}

% Co-Training / Daten
\St{old}{12.2}{\yC}{above}{Render v1}{Greifpunkt verworfen}
\St{part}{13.8}{-5.65}{right}{Cosmos-Transfer}{\Detail{nie Ende-zu-Ende}}
\St{fail}{15.1}{\yC}{above}{Synth. Datensatz v1}{kein Co-Training}
\St{part}{17.4}{\yC}{below}{Externe Synthesedaten}{Fingerzuordnung vertauscht}
\St{ok}{19.7}{\yC}{above}{Render-Daten v3}{187 Ep. (L72)}
\Mk{ok}{21.2,\yT}\Mk{ok}{21.2,\yN}\Mk{ok}{21.2,\yC}
\node[lbl, below=3.2mm] at (21.2,\yC) {{\bfseries Co-Training v3}\Detail{\\4× 44k (L73–76)}};
\ifkurz % Folie: größere Schrift, kürzerer Titel, sonst stößt er an den Kasten
\St*{ok}{22.7}{\yC}{above}{Lücke halbiert}{MSE 0,024→0,012}
\else
\St*{ok}{22.7}{\yC}{above}{Bildlücke halbiert}{MSE 0,024→0,012}
\fi

% Simulation / Szene
\St!{ok}{1.6}{\yS}{above}{G1 + DEX3}{in Isaac Lab}
\node[lbl, lSim, font=\fsLbl\itshape] at (6.95,-8.07) {Fehlannahme: Finger-Vorzeichen – Hand schließt nie};
\St{ok}{8.2}{\yS}{above}{Würfel-Startlage}{Env-Versatz behoben (L8/9)}
\St*{ok}{9.6}{\yS}{below}{Kameras verdreht}{96–137° (L10–18)}
\St{part}{11.0}{\yS}{above}{Greifphysik}{hebt 7,9 cm (L29)}
\St*{ok}{13.95}{\yS}{above}{Geometrieabgleich}{Finger, Becken, Basis}
\St*{ok}{18.3}{\yS}{above}{Würfel 3 cm zu weit innen}{Demo hebt 6/20 (L62/69)}
\St{part}{20.0}{\yS}{below}{Szenenabgleich}{Handfehler −40\,\% (L70)}

% Evaluation / Diagnose
\St!{fail}{1.8}{\yD}{above}{Erste Sim-Eval}{Modell T1: 0/20}
\St{old}{3.95}{-10.95}{left}{SigLIP-Gap}{0,26 / 0,43}
\node[lbl, lDiag, font=\fsLbl\itshape] at (11.3,-11.22) {Fehlannahme: SigLIP-Abstand als Maß der Lücke (0,22 sah »gut« aus)};
\St*{ok}{6.0}{\yD}{above}{RoboCasa-Referenz}{47,7\,\% ≈ 47,8\,\%}
\St{fail}{8.4}{\yD}{above}{Messpunkt 3× falsch}{(L24–28)}
\St*{ok}{9.8}{\yD}{below}{Fingerspannen-Gate}{echt 100\,\%, Sim 19\,\% (L32)}
\St{part}{11.3}{\yD}{above}{T3 im Closed Loop}{0/10 (L34)}
\ifkurz % Folie: unter die Linie, sonst stößt sie an »Bild-Ablation«
\St*{fail}{16.9}{\yD}{below}{Demo-Replay}{}
\else
\St*{fail}{16.9}{\yD}{above}{Demo-Replay}{hebt 1/16 (L57)}
\fi
\St*{ok}{19.0}{\yD}{above}{Bild-Ablation}{Fehler ×13 / ×36}
\St{ok}{20.6}{\yD}{above}{Fester Seed}{Gate (L80)}
\St*{ok}{21.85}{\yD}{below}{Vorregistrierter Test}{p = 5·10⁻⁶ (L96)}
\St{part}{23.1}{\yD}{above}{60 Testlagen}{2/60 (L97)}

% RL
\St!{ok}{5.0}{\yR}{above}{RL-Plan (FPO)}{12.06.}
\St{part}{8.0}{\yR}{below}{RL-Pipeline läuft}{ohne Lerneffekt}
\node[lbl, lRL!80!black, font=\fsLbl\itshape, above=1.2mm] at (18.5,\yR)
  {wartet auf ein Griffsignal (erstmals 2/60)};

% Doku
\St{ok}{0.35}{\yK}{above}{README}{}
\St{ok}{12.6}{\yK}{above}{Doku neu geordnet}{18.08.}
\St{ok}{17.6}{\yK}{above}{Ausarbeitung neu}{11.10.}
\St!{ok}{22.6}{\yK}{above}{Abgabe-Review}{30.10.}

% ════════════════════════════ STAND HEUTE ═════════════════════════════
\begin{scope}[shift={(\xH,0)}]
  \draw[black, line width=1.2pt, fill=hubbg, rounded corners=3mm] (0,1.25) rectangle (2.75,-14.05);
  \node[font=\fontsize{9}{10}\selectfont\sffamily\bfseries, anchor=north, align=center]
    at (1.375,1.1) {Stand\\[-1pt]01.11.2026};
  % Einmündungen
  \foreach \y/\cc in {\yI/lInfra,\yT/lTrain,\yN/lNsieben,\yC/lCotr,\yS/lSim,\yD/lDiag,\yK/lDoku}
    \fill[\cc] (0,\y) circle (1.3mm);
  \node[font=\fsLbl, text width=24.5mm, align=left, anchor=north west] at (0.15,-0.35) {%
    \ifkurz {\bfseries Bestes Modell}\\N1.6 @44k\\[3pt]{\bfseries Sim-Kandidat}\\co-trainiert
    \else {\bfseries Bestes Modell}\\N1.6 vision-v2 (T3) @44k\\Val-MSE 0,00482\\[2pt]
    {\bfseries Bester Sim-Kandidat}\\N1.6 co-trainiert,\\cube4 @44k (L74)\fi};
  % Trichter Closed Loop (L97, 60 Testlagen)
  \node[font=\fsHub\bfseries, anchor=west] at (0.18,-3.95) {Closed Loop (L97)};
  \foreach \i/\lab/\k/\n in {0/Kontakt/60/60, 1/{Hand ≤ 4 cm}/58/60, 2/gehoben/6/60,
      3/frei gehoben/2/60, 4/gestapelt/0/60} {
    \pgfmathsetmacro\w{2.35*\k/\n}
    \fill[black!10] (0.2,{-4.45-\i*0.62}) rectangle (2.55,{-4.68-\i*0.62});
    \ifnum\k>0 \fill[lDiag] (0.2,{-4.45-\i*0.62}) rectangle ({0.2+\w},{-4.68-\i*0.62}); \fi
    \node[font=\fsHub, anchor=south west, inner sep=0.5pt] at (0.18,{-4.43-\i*0.62}) {\lab};
    \node[font=\fsHub\bfseries, anchor=south east, inner sep=0.5pt] at (2.57,{-4.43-\i*0.62}) {\k/\n};
  }
  \node[font=\fsLbl, text width=24.5mm, align=left, anchor=north west] at (0.15,-7.65) {%
    \ifkurz
    {\bfseries Blocker}\\
    {\color{sPart!70!black}\textbullet} Würfel versetzt\\
    {\color{sPart!70!black}\textbullet} Bild-Lücke\\[4pt]
    {\bfseries Offen}\\
    {\color{sFail}\textbullet} Stapeln: 0\\
    {\color{sFail}\textbullet} kein Robotertest
    \else
    {\bfseries Zwei Blocker}\\
    {\color{sPart!70!black}\textbullet} Würfel 3 cm versetzt\\
    \hphantom{\textbullet} Demo hebt 6/20 (L69)\\
    {\color{sPart!70!black}\textbullet} Bild-Lücke ×13\\
    \hphantom{\textbullet} halbiert, nicht weg\\[3pt]
    {\bfseries Offen}\\
    {\color{sFail}\textbullet} Griff 2/60, Stapeln 0\\
    {\color{sFail}\textbullet} kein Robotertest\\
    {\color{black!50}\textbullet} RL ohne Ergebnis\\[3pt]
    {\bfseries Nächster Schritt}\\
    mehr Render-Episoden,\\dann Test am G1
    \fi};
\end{scope}

\end{tikzpicture}%
